\section*{Bab \ref{ch:b2:neurons-synapses} --- Neuron, Potensial Aksi, dan Sinapsis}

\begin{solution}{exo:b2:neurons-synapses:1}
\omterm{def:b2:neurons-synapses:neuron}{Dendritnya} menerima \omterm{def:b2:neurons-synapses:synapse}{sinapsis}; somanya menahan nukleusnya dan
memadukan; \omterm{prop:b2:neurons-synapses:integration}{bukit aksonnya} memutuskan; \omterm{def:b2:neurons-synapses:neuron}{aksonnya} menghantar;
terminalnya melepaskan penghantar. \omterm{def:b2:neurons-synapses:neuron}{Mielin} adalah membran terlilit
sebuah sel \omterm{def:b2:neurons-synapses:neuron}{glia}, yang nyaris bebas sitoplasma; ia menyekat \omterm{def:b2:neurons-synapses:neuron}{aksonnya}
sehingga impulsnya meloncat di antara nodusnya dan berjalan seratus
kali lebih cepat dengan ongkos yang lebih rendah.
\end{solution}

\begin{solution}{exo:b2:neurons-synapses:2}
Depolarisasi sampai ambangnya; fase naiknya --- kanal natrium
bergerbang tegangan membuka, umpan balik positif; puncaknya mendekati
$E_{\mathrm{Na}}$ --- kanal natriumnya menonaktifkan diri; fase
turunnya --- kanal kalium bergerbang tegangan membuka; ayunan bawahnya
--- kanal kaliumnya masih terbuka, dekat $E_{\mathrm K}$; pemulihannya
ketika kanalnya menutup. Semua-atau-tidak-sama-sekali karena umpan
balik positifnya entah berjalan sampai tuntas atau tidak sama sekali;
satu arah karena membran yang baru dilewatinya bersifat refrakter
(kanal natriumnya nonaktif).
\end{solution}

\begin{solution}{exo:b2:neurons-synapses:3}
Impulsnya mencapai terminalnya; kanal kalsium bergerbang tegangan
membuka; kalsiumnya memicu peleburan vesikelnya; penghantarnya
berdifusi menyeberangi celahnya; mengikat reseptor pascasinapsisnya;
kanalnya membuka (atau sebuah \omterm{prop:b2:cell-signalling:gpcr}{protein G} menyakelar); \omterm{def:b2:neurons-synapses:synapse}{potensial
pascasinapsisnya}; penghantarnya disingkirkan enzim atau pengangkut;
membran vesikelnya diambil kembali.
\end{solution}

\begin{solution}{exo:b2:neurons-synapses:4}
Listrik: tanpa tundaan, biasanya dua arah, tanpa penguatan, tanpa
penghambatan (hanya arus yang lewat). Kimia: setengah milidetik, satu
arah, dengan penguatan (satu impuls melepaskan ratusan kuantum), dengan
penghambatan yang mungkin (kanal bagi klorida atau kalium), dan dapat
diubah.
\end{solution}

\begin{solution}{exo:b2:neurons-synapses:5}
Dari $E = (61.5/z)\log_{10}(c_{\text{luar}}/c_{\text{dalam}})$: kalium
\qty{-89}{mV}; natrium \qty{+61}{mV}; klorida, dengan $z =
-1$,
$-61.5\log_{10} 11 = \qty{-64}{mV}$; kalsium, dengan $z = 2$,
$30.75\times 4.3 = \qty{+132}{mV}$. Pada \qty{-70}{mV}: kaliumnya
keluar, natriumnya masuk, kloridanya keluar sedikit (membrannya berada
di bawah kesetimbangannya), kalsiumnya masuk dengan kuat.
\end{solution}

\begin{solution}{exo:b2:neurons-synapses:6}
$20 : 1$: $(20\times(-89) + 61)/21 = \qty{-82}{mV}$, yaitu keadaan
istirahatnya dekat $E_{\mathrm K}$. $1 : 20$: $(-89 + 20\times 61)/21 =
\qty{+54}{mV}$, yaitu puncaknya, dekat $E_{\mathrm{Na}}$. $50 : 1$: $(50\times
(-89) + 61)/51 = \qty{-86}{mV}$, yaitu ayunan bawahnya dengan kanal
kalium tambahan yang terbuka.
\end{solution}

\begin{solution}{exo:b2:neurons-synapses:7}
$\sqrt{500} = 22$: \qty{22}{m/s}. Untuk mencapai \qty{100}{m/s} sebuah
\omterm{def:b2:neurons-synapses:neuron}{akson} tanpa \omterm{def:b2:neurons-synapses:neuron}{mielin} memerlukan $d = 10^{4}$ \unit{\micro m}, yaitu satu
sentimeter. Satu meter memerlukan \qty{10}{ms} pada serabut
bermielinnya, \qty{45}{ms} pada \omterm{def:b2:neurons-synapses:neuron}{akson} cumi-cuminya, dan satu detik
penuh pada sebuah serabut \qty{1}{\micro m}.
\end{solution}

\begin{solution}{exo:b2:neurons-synapses:8}
$m = \ln(200/74) = 1.0$; $P(1) = e^{-1} = 0.37$, $P(2) = 0.18$, $P(3)
= 0.06$. Tanggapan reratanya $1.0\times 0.5 = \qty{0.5}{mV}$.
\end{solution}

\begin{solution}{exo:b2:neurons-synapses:9}
$1/\qty{2}{ms} = 500$ impuls tiap detik. Kekuatan sentuhannya
disandikan oleh frekuensi impulsnya (dan oleh berapa banyak \omterm{def:b2:neurons-synapses:neuron}{neuron} yang
direkrut); impulsnya sendiri serupa.
\end{solution}

\begin{solution}{exo:b2:neurons-synapses:10}
$\lambda = \sqrt{R_m d/4R_i}$: \qty{0.5}{mm}, \qty{1.6}{mm},
\qty{11}{mm}. Dengan \omterm{def:b2:neurons-synapses:neuron}{mielin}, \qty{10}{\micro m}: $1.6\times\sqrt{200} =
\qty{22}{mm}$. Sebuah antarnodus \qty{1}{mm} hanya kehilangan $1 - e^{-1/22}
= \qty{4}{\%}$ isyaratnya, sehingga nodus berikutnya dengan mudah
dibawa sampai ambangnya. Melintasi \qty{3}{mm} tanpa \omterm{def:b2:neurons-synapses:neuron}{mielin} isyaratnya
turun menjadi $e^{-3/1.6} = \qty{15}{\%}$ --- dari sebuah impuls
\qty{100}{mV}, yaitu kira-kira ambang \qty{15}{mV}: penghantarannya
gagal atau menjadi tak terandalkan.
\end{solution}

\begin{solution}{exo:b2:neurons-synapses:11}
Tetrodotoksin: tanpa arus natrium, tanpa impuls, tanpa penghantaran.
Tetraetilamonium: tanpa arus kalium, impulsnya memanjang, lebih banyak
penghantar dilepaskan tiap impulsnya. Tanpa kalsium: impulsnya normal,
tanpa pelepasan, tanpa penghantaran. Kurare: pelepasannya normal,
reseptornya terhalang, tanpa potensial lempeng-ujung --- kelumpuhan.
Penghambat kolinesterase: asetilkolinnya bertahan, reseptornya tetap
terbuka, ototnya berdepolarisasi lalu tidak dapat berepolarisasi ---
kelumpuhan oleh kelebihan. Toksin botulinum: vesikelnya tidak dapat
melebur, tanpa pelepasan --- kelumpuhan.
\end{solution}

\begin{solution}{exo:b2:neurons-synapses:12}
Tundaannya membeli: penghantaran hanya dalam satu arah; penguatan
(sebuah impuls tunggal melepaskan ratusan kuantum dan dapat
menggerakkan sel yang lebih besar); penghambatan, yang tidak dapat
dikerjakan sebuah persambungan listrik; dan keplastisan, karena jumlah
kuantum dan reseptornya dapat berubah menurut pemakaiannya, dan itulah
cara untaiannya belajar. \omterm{def:b2:neurons-synapses:synapse}{Sinapsis} listrik dipakai di tempat
keserentakan dan kecepatan lebih berarti daripada perhitungan: untaian
lolos, otot \omterm{def:b2:heart:anatomy}{jantungnya}, beberapa jaringan penghambat.
\end{solution}

\begin{solution}{pb:b2:neurons-synapses:1}
\textbf{1.} $E_{\mathrm K} = 61.5\log_{10}(4/140) = \qty{-95}{mV}$;
$E_{\mathrm{Na}} = 61.5\log_{10}(145/12) = \qty{+67}{mV}$.
\textbf{2.} $(25\times(-95) + 67)/26 = \qty{-89}{mV}$.
\textbf{3.} $E_{\mathrm K} = 61.5\log_{10}(8/140) = \qty{-76}{mV}$;
$V = (25\times(-76) + 67)/26 = \qty{-71}{mV}$: lebih dekat ke
ambangnya, sehingga mula-mula lebih terangsang --- dan, jika ia
bertahan, kurang terangsang, karena kanal natriumnya nonaktif pada aras
yang terdepolarisasi itu.
\textbf{4.} Landaiannya habis selama berjam-jam ketika natriumnya
merembes masuk dan kaliumnya keluar; potensialnya menghanyut menuju
nol; dengan pompanya berhenti, anion selnya yang tak telap menarik
masuk klorida dan air, dan selnya menggembung.
\textbf{5.} Permukaannya tiap milimeter $\pi\times 15\times 10^{-6}\times
10^{-3} = \qty{4.7e-8}{m^2}$; $C = \qty{4.7e-10}{F}$; $Q = CV =
\qty{4.7e-11}{C}$: $2.9\times 10^{8}$ ion.
\textbf{6.} Volumenya tiap milimeter $\pi(7.5\times 10^{-6})^{2}\times
10^{-3} = \qty{1.8e-13}{m^3} = \qty{1.8e-10}{L}$, yang menahan $12\times
10^{-3}\times 1.8\times 10^{-10} = 2.1\times 10^{-12}$ mol, yaitu $1.3\times
10^{12}$ ion: sebuah perubahan sebesar $2\times 10^{-4}$.
\textbf{7.} $2.9\times 10^{8}/3 \approx 10^{8}$ ATP tiap milimeter tiap
impulsnya.
\textbf{8.} $0.9/90 = \qty{10}{ms}$ tiap arahnya.
\textbf{9.} Telanjang: $\sqrt{1\times 15\times 10^{-6}/4} = \qty{1.9}{mm}$;
bermielin: $\times\sqrt{250} = \qty{31}{mm}$.
\textbf{10.} Telanjang: $e^{-1.5/1.9} = 0.45$, \qty{45}{mV};
bermielin: $e^{-1.5/31} = 0.95$, \qty{95}{mV}. Keduanya melampaui
ambang \qty{15}{mV}; tetapi \omterm{def:b2:neurons-synapses:neuron}{akson} telanjangnya harus memuati seluruh
membrannya di sepanjang jalannya, dan itulah yang membuatnya lambat,
sedangkan di bawah \omterm{def:b2:neurons-synapses:neuron}{mielinnya} nyaris tidak ada yang hilang atau dimuati.
\textbf{11.} $\sqrt{15} = \qty{3.9}{m/s}$: \qty{0.23}{s} tiap arahnya,
hampir setengah detik bagi perjalanan bolak-baliknya --- terlalu lambat
bagi sebuah refleks yang harus menangkap sebuah tersandung.
\textbf{12.} $e^{-4/1.9} = 0.12$: \qty{12}{mV} pada nodus yang jauh, di
bawah ambangnya --- impulsnya terhalang dan refleksnya hilang.
\textbf{13.} Di belakang impulsnya kanal natriumnya nonaktif selama
satu atau dua milidetik, sehingga membran yang baru dilewatinya tidak
dapat dirangsang ulang dan gelombangnya hanya dapat bergerak maju.
\textbf{14.} $m = \ln(300/111) = 1.0$.
\textbf{15.} $P(1) = 0.37$, $P(2) = 0.18$, $P(3) = 0.06$: kira-kira
110, 55, dan 18 dari 300 percobaannya.
\textbf{16.} $1.0\times 0.4 = \qty{0.4}{mV}$.
\textbf{17.} $P(0) = e^{-60} \approx 10^{-26}$: tidak pernah;
tanggapan reratanya \qty{24}{mV}, di atas ambang \qty{15}{mV}: satu
\omterm{def:b2:neurons-synapses:synapse}{sinapsis} semacam itu memicu \omterm{def:b2:neurons-synapses:neuron}{neuron} motoriknya.
\textbf{18.} Dua puluh \omterm{def:b2:neurons-synapses:synapse}{sinapsis} yang aktif bersama menjumlahkan arusnya
(penjumlahan ruang), secara kurang lurus ketika potensialnya mendekati
potensial balikan kanalnya; bahkan pada $m = 1$ masing-masingnya akan
memberikan \qty{8}{mV}, dan pada $m$ yang normal jauh lebih banyak
daripada ambangnya.
\textbf{19.} Membuka kanal klorida pada $E_{\mathrm{Cl}} = V_{\text{istirahat}}$
menambahkan konduktans tanpa memindahkan potensialnya; arus
perangsangnya kini mengalir lewat sebuah membran yang lebih bocor lalu
menghasilkan depolarisasi yang lebih kecil --- penghambatan pirau.
\textbf{20.} $10 + 0.6 + 10 + 0.6 + 5 = \qty{26}{ms}$, berbanding
\qty{30}{ms} yang terukur.
\textbf{21.} Pengaktifan reseptornya sendiri di dalam kumparan ototnya,
perambatan \omterm{prop:b2:neurons-synapses:ap}{potensial aksi} \omterm{def:b2:cell-differentiation:myogenesis}{serabut ototnya} di sepanjang serabutnya
(beberapa meter tiap detik sepanjang beberapa sentimeter), dan tundaan
mekanis sebelum gayanya muncul.
\textbf{22.} Tiap \omterm{def:b2:neurons-synapses:synapse}{sinapsis} menambahkan setengah milidetik dan sebuah
interneuron menambahkan sebuah sel; tetapi interneuron membolehkan
isyaratnya dibalik (dengan menghambat lawannya), dipadukan dengan yang
lain, dan diselia otaknya --- sebuah refleks yang dapat dibentuk.
\textbf{23.} $90 - 6.2 = \qty{84}{ms}$ penghantaran bagi \qty{1.8}{m}:
kira-kira \qty{21}{m/s}.
\textbf{24.} Separuh arus masuknya pada tiap tegangannya: ambangnya
naik, impulsnya lebih kecil dan lebih lambat, faktor keamanannya pada
tiap nodus turun, dan impulsnya mungkin gagal; refleksnya melemah atau
lenyap.
\textbf{25.} \omterm{thm:b2:neurons-synapses:chord}{Potensial istirahatnya} \qty{-89}{mV}; \qty{10}{ms} tiap
arahnya; $m = 1.0$; tundaan yang dihitung \qty{26}{ms}.
\end{solution}
